\documentclass{article}
\usepackage{amssymb}
\usepackage{pifont}
\usepackage{booktabs}
\usepackage{subcaption}
\usepackage{colortbl}
\usepackage[dvipsnames]{xcolor}
\usepackage{graphicx}
\usepackage{siunitx}
\captionsetup[subfigure]{position=bottom}
\begin{document}
\newcommand{\haliver}{HaliVer}
\begin{table}[t]
\centering
\newcommand{\widthPadre}{0.9}
\subfloat[\label{tab:blaslevel1}\texttt{CLBLAS level 1 kernels}]{
\resizebox{\widthPadre\textwidth}{!}{
\begin{tabular}{l|lll|l|rr|rr S[table-format=-4.1]|rr S[table-format=-4.1]|rr S[table-format=-4.1]|rr S[table-format=-4.1]r}
\hline
 & & & & & \multicolumn{2}{c|}{\textbf{Base}}& \multicolumn{3}{c}{\textbf{Unique}}& \multicolumn{3}{c}{\textbf{Immutable}}& \multicolumn{3}{c}{\textbf{Extract}}& \multicolumn{3}{c}{\textbf{All}}\\
\textbf{Kernel} & $\#_{imm}$ & $\#_{p}$ & $\#_{A}$ & \textbf{Result} & \textbf{\#} & \textbf{T} & \textbf{\#} & \textbf{T} & \textbf{Speedup} & \textbf{\#} & \textbf{T} & \textbf{Speedup} & \textbf{\#} & \textbf{T} & \textbf{Speedup} & \textbf{\#} & \textbf{T} & \textbf{Speedup}\\
\hline
xamax &  1 &  4 &  31 & \checkmark&  10 & 34&  10&  30& 1.1&  10&  29& 1.2&  10& \cellcolor{ForestGreen!25} 24& 1.4&  10& \cellcolor{ForestGreen!25} 24& 1.4 \\
\hline
xasum &  1 &  2 &  37 & \checkmark&  10 & 20&  10&  19& 1.1&  10&  16& 1.2&  10& \cellcolor{ForestGreen!25} 13& 1.5&  10& \cellcolor{ForestGreen!25} 13& 1.5 \\
\hline
xaxpy$^\dagger$ &  1 &  1 &  16 & \checkmark&  9 & 66&  10&  19& 3.3&  10&  24& 2.6&  9&  13& 4.8&  10& \cellcolor{ForestGreen!25} 9& 6.9 \\
\hline
 & & && $\times$& 1 & 23& 0 &  && 0 &  && 1 & 13 && 0 &  & \\
\hline
xcopy$^\dagger$ &  1 &  1 &  15 & \checkmark&  8 & 52&  10&  16& 6.3&  10&  15& 6.7&  10&  22& 4.6&  10& \cellcolor{ForestGreen!25} 10& 10.1 \\
\hline
 & & && $\times$& 2 & 296& 0 &  && 0 &  && 0 &  && 0 &  & \\
\hline
xdot &  2 &  2 &  24 & \checkmark&  10 & 20&  10&  18& 1.1&  10&  16& 1.2&  10&  15& 1.3&  10& \cellcolor{ForestGreen!25} 12& 1.7 \\
\hline
xhad$^\dagger$ &  2 &  1 &  19 & \checkmark&  8 & 450&  10&  73& 8.0&  10&  48& 12.1&  10&  54& 10.8&  10& \cellcolor{ForestGreen!25} 12& 48.4 \\
\hline
 & & && $\times$& 2 & 1104& 0 &  && 0 &  && 0 &  && 0 &  & \\
\hline
xnrm2 &  1 &  2 &  54 & \checkmark&  10 & 21&  10&  20& 1.1&  10&  18& 1.2&  10&  15& 1.4&  10& \cellcolor{ForestGreen!25} 14& 1.5 \\
\hline
xswap$^\dagger$ &  0 &  2 &  19 & \checkmark&  10 & 66&  10&  32& 2.1& - & - & &  10&  28& 2.4&  10& \cellcolor{ForestGreen!25} 12& 5.5 \\
\hline
xscal$^\dagger$ &  0 &  1 &  13 & \checkmark&  10 & 8& - & - & & - & - & &  10& \cellcolor{ForestGreen!25}\cellcolor{BrickRed!25} 10& 0.8& - & - &  \\
\hline
\hline
Total & & & & & & 913 & & 227 & 4.0 & & 166 & 5.1 & & 194 & 4.7 & & 106 & 8.5  \\
\hline
\end{tabular}
}
}
\\
\subfloat[\label{tab:blaslevel2}\texttt{CLBLAS level 2 kernels}]{
\resizebox{\widthPadre\textwidth}{!}{
\begin{tabular}{l|lll|l|rr|rr S[table-format=-4.1]|rr S[table-format=-4.1]|rr S[table-format=-4.1]|rr S[table-format=-4.1]r}
\hline
 & & & & & \multicolumn{2}{c|}{\textbf{Base}}& \multicolumn{3}{c}{\textbf{Unique}}& \multicolumn{3}{c}{\textbf{Immutable}}& \multicolumn{3}{c}{\textbf{Extract}}& \multicolumn{3}{c}{\textbf{All}}\\
\textbf{Kernel} & $\#_{imm}$ & $\#_{p}$ & $\#_{A}$ & \textbf{Result} & \textbf{\#} & \textbf{T} & \textbf{\#} & \textbf{T} & \textbf{Speedup} & \textbf{\#} & \textbf{T} & \textbf{Speedup} & \textbf{\#} & \textbf{T} & \textbf{Speedup} & \textbf{\#} & \textbf{T} & \textbf{Speedup}\\
\hline
xger$^\dagger$ &  2 &  3 &  69 & \checkmark& \cellcolor{BrickRed!25} 0 & & \cellcolor{ForestGreen!25} 9&  369& 5.3& \cellcolor{ForestGreen!25} 5&  981& 2.4& \cellcolor{ForestGreen!25} 3&  1443& 1.7& \cellcolor{ForestGreen!25} 10& \cellcolor{ForestGreen!25} 14& 151.5 \\
\hline
 & & && $\times$& 8 & 1751& 1 & 701 && 5 & 803 && 7 & 1149 && 0 &  & \\
\hline
 & & && T.O.& 2 & -& 0 &  && 0 &  && 0 &  && 0 &  & \\
\hline
xher$^\dagger$ &  1 &  3 &  37 & \checkmark&  3 & 1507&  10&  261& 3.0&  9&  280& 2.6&  5&  401& 2.2&  10& \cellcolor{ForestGreen!25} 42& 18.5 \\
\hline
 & & && $\times$& 7 & 465& 0 &  && 1 & 495 && 5 & 315 && 0 &  & \\
\hline
xher2$^\dagger$ &  2 &  5 &  56 & \checkmark& \cellcolor{BrickRed!25} 0 & & \cellcolor{ForestGreen!25} 2&  922& 0.7& \cellcolor{ForestGreen!25} 4&  731& 1.4& \cellcolor{BrickRed!25} 0&  & 1.4& \cellcolor{ForestGreen!25} 10& \cellcolor{ForestGreen!25} 48& 21.5 \\
\hline
 & & && $\times$& 10 & 1030& 8 & 1638 && 6 & 739 && 10 & 733 && 0 &  & \\
\hline
xtrsv$^\dagger$ &  2 &  3 &  68 & \checkmark&  8 & 1381&  10& \cellcolor{ForestGreen!25} 330& 5.5&  10&  610& 3.0&  9&  1217& 1.4&  10&  455& 4.0 \\
\hline
 & & && $\times$& 0 & & 0 &  && 0 &  && 1 & 1933 && 0 &  & \\
\hline
 & & && T.O.& 2 & -& 0 &  && 0 &  && 0 &  && 0 &  & \\
\hline
xgemv &  2 &  3 &  80 & \checkmark&  6 & 1664&  10&  380& 4.6&  10&  159& 11.1&  9&  787& 2.1&  10& \cellcolor{ForestGreen!25} 135& 13.0 \\
\hline
 & & && $\times$& 4 & 1905& 0 &  && 0 &  && 1 & 1411 && 0 &  & \\
\hline
\hline
Total & & & & & & 7514 & & 2867 & 2.6 & & 2699 & 2.8 & & 4466 & 1.7 & & 694 & 10.8  \\
\hline
\end{tabular}
}
}
\\
\caption{\label{tab:blas-results}Verification results for the clBLAS GPU kernels for level 1 and 2. 
\textbf{Base} is the baseline configuration without any extra features. It is compared against versions with unique type qualifiers (\textbf{Unique}), immutable arrays (\textbf{Immutable}), extracted kernel bodies (\textbf{Extract}), and a configuration enabling all features simultaneously (\textbf{All}).
The $\dagger$ indicates we (try to) prove functional correctness in addition to memory safety.
% A name marked with a $\dagger$ indicates that \textsc{VerCors} attempts to prove both memory safety and functional correctness of the kernel; otherwise, it only attempts to prove memory safety.
% $T_t$ denotes the average total verification time of VerCors, and $T_v$ the average time spent in the VerCors backend. Speedup$_v$ is the speedup of backend time relative to the baseline.
$\#_{imm}$ and $\#_p$ is the number of immutable and normal pointer arrays, respectively. $\#_A$ the is the total number of annotations.}
\end{table}
\begin{figure}[t]
\centering
\end{figure}
\begin{figure}[t]
\centering
\subfloat[\label{fig:blas-barplot-level1}Level 1]{\includegraphics[width=0.95\textwidth]{/home/lars/data/Improved-Deductive-Verification-of-Parallel-Programs---CAV2026/experiments/ClBlast/results/blas_level1_barplot.pdf}}\\
\subfloat[\label{fig:blas-barplot-level2}Level 2]{\includegraphics[width=0.95\textwidth]{/home/lars/data/Improved-Deductive-Verification-of-Parallel-Programs---CAV2026/experiments/ClBlast/results/blas_level2_barplot.pdf}}
\caption{\label{fig:blas-results}Verification backend times for the clBLAS GPU kernels for level 1 and 2, normalized so Base=100\%. \textbf{Base} is the baseline configuration without any extra features. It is compared against versions with unique type qualifiers (\textbf{Unique}), immutable arrays (\textbf{Immutable}), extracted kernel bodies (\textbf{Extract}), and a configuration enabling all features simultaneously (\textbf{All}). Bars show mean with standard error of the mean. Below the bars: \checkmark = 10/10 successful, \ding{55}  = 0/10 successful, else number of successful runs out of 10. Next to base bars: average time in seconds of the 10 base runs. The $\dagger$ indicates we (try to) prove functional correctness in addition to memory safety.}
\end{figure}
\end{document}